Informatics Engineering graduate with a strong foundation in data-driven problem-solving and AI implementation. Proven track record in bridging technical execution with business requirements to deliver impactful enterprise solutions. Experienced in leading cross-functional teams, driving operational efficiency, and managing comprehensive tech community initiatives. Passionate about leveraging emerging technologies, data analysis, and UI/UX prototyping to optimize business workflows and support strategic decision-making

\section*{Education}
\resumeSubheading{Bachelor's Degree in Informatics Engineering | Institut Teknologi Sumatera}{Aug 2022 - May 2026}{Lampung, Indonesia}{}
\begin{itemize}
  \item GPA: 3.63/4.00
  \item Focus: Deep Learning, Machine Learning, NLP, and Computer Vision
  % \item Relevant Coursework: Artificial Intelligence, Machine Learning, Deep Learning, Natural Language Processing, Multimodal Machine Learning, Data Visualization and Information
\end{itemize}

\section*{Experience}
\resumeSubheading{PT Pertamina Hulu Rokan - AI Engineer Intern }{Jun 2025 - Jul 2025}{Pekanbaru, Indonesia}{}
\begin{itemize}
  \item Spearheaded the evaluation of Retrieval-Augmented Generation (RAG) pipelines for enterprise document management, identifying and mitigating AI hallucination risks to ensure data reliability
  \item Facilitated cross-functional discussions between technical mentors and business stakeholders to align system capabilities with user requirements, successfully delivering alternative solutions that met operational needs
  \item Optimized document QA workflows, achieving a 73.33\% accuracy rate while maintaining clear communication regarding system limitations and risk mitigation strategies to upper management
\end{itemize}

% \section*{Teaching Experience}
\resumeSubheading{UPA TIK ITERA - Teaching Assistant}{Aug 2023 - Present}{Lampung, Indonesia}{}
\begin{itemize}
  \item \textbf{Operating Systems}\\
  Facilitated hands-on laboratory sessions focusing on process management, memory allocation, and operating system fundamentals

  \item \textbf{Introduction to Computing \& Software I (PKS I)}\\
  Guided students in procedural programming, computational thinking, and fundamental programming logic\\
  \textit{Credential ID: B/5359/IT9.E2/TA.00.00/2024}

  \item \textbf{Interaction Design}\\
  Reviewed student assignments and provided structured feedback on clarity, and project quality\\
  \textit{Credential ID: 33133/IT9.3.3/TA.00.00/2025}
\end{itemize}

\section*{Projects}
\resumeProject{Multimodal Deep Learning for Self-Harm Meme Detection $|$ \href{https://github.com/elsaelisa09/multimodal-meme-selfharm}{\footnotesize [GitHub]}}{Aug 2025 - Apr 2026}{Tech: CLIP, ELECTRA, Gradio}
\vspace{-0.3em}
\begin{itemize}
  % \item Built and curated a custom dataset of 2,024 memes by collecting and cleaning image-text samples, then performing LLM-assisted annotation for self-harm and non-self-harm labels based on predefined criteria
  % \item Developed a multimodal binary classification system for a final thesis project to identify self-harm versus non-self-harm memes using combined visual and textual representations
  % % \item Engineered an intermediate fusion architecture via concatenation, leveraging CLIP ViT-B/32 for image encoding and ELECTRA-base for text encoding
  % \item Trained and evaluated the model on a stratified dataset of 2,024 memes, achieving an outstanding 96.1\% F1-score and successfully outperforming both single-modality baselines (CLIP at 91.5\% and ELECTRA at 92.0\%)
  % % \item Deployed an interactive web-based inference demonstration of the model utilizing Gradio on Hugging Face Spaces
  \item Built and curated a custom multimodal dataset by collecting, cleaning, and annotating image-text samples using LLM-assisted labeling based on predefined criteria
  \item Reviewed and validated annotation results to ensure labeling consistency and data reliability
  \item Conducted model evaluation and error analysis across multimodal and single-modality baselines, achieving 96.1\% F1-score
\end{itemize}

\resumeProject{Multimodal City Branding Analysis $|$ \href{https://github.com/elsaelisa09/city-branding-multimodal}{\footnotesize [GitHub]}}{May 2026 - Present}{Tech: Image-Text Representation Learning, Social Media Analysis}
\vspace{-0.3em}
\begin{itemize}
  \item Developed a multimodal deep learning model to analyze South Lampung tourism image and caption data from Instagram
  % \item Used openai/clip-vit-base-patch32 for image representation and cardiffnlp/twitter-xlm-roberta-base for text representation
  % \item Handled class imbalance using Focal Loss with class weights, especially due to the dominance of the Attraction class compared to other categories
  \item Achieved approximately 89.2\% validation accuracy at the selected best epoch based on validation loss
  \item Collaborated with researchers and students from Urban and Regional Planning and Informatics Engineering
\end{itemize}

\resumeProject{Student Face Recognition with ArcFace and SVM $|$ \href{https://github.com/elsaelisa09/student-face-recognition-arcface-svm}{\footnotesize [GitHub]}}{Jun 2026}{Tech: InsightFace ArcFace, Scikit-learn, OpenCV, SVM, Jupyter Notebook}
\vspace{-0.3em}
\begin{itemize}
  \item Developed an end-to-end facial identification system to automate the recognition of 70 student identities from uploaded facial images
  \item Engineered a pipeline using offline image augmentation, 512-dimensional InsightFace ArcFace Buffalo\_L embeddings, and an RBF-kernel SVM classifier with leakage-resistant Stratified Group K-Fold evaluation
  \item Processed 1,946 facial samples and achieved 82.0\% accuracy, 85.4\% weighted precision, and 82.4\% weighted F1-score
\end{itemize}

\section*{Leadership \& Organizations}
\resumeSubheading{Head of Academic \& Scholarship Division — HMIF ITERA}{Jan 2025 – Dec 2025}{Lampung, Indonesia}{}
\vspace{-0.3em}
\begin{itemize}
  \item Led 4 academic and scholarship-related programs, served as a speaker in a scholarship-sharing session, and managed mentoring initiatives that contributed to 2 students receiving the Pertamina Sobat Bumi Scholarship
\end{itemize}

\resumeSubheading{Secretary of Desa Energi Berdikari - Sobat Bumi Lampung}{Nov 2024 - Dec 2025}{Lampung, Indonesia}{}
\vspace{-0.3em}
\begin{itemize}
  \item Directed cross-functional coordination among internal divisions to streamline project planning and finalize comprehensive proposals for the Desa Energi Berdikari (DEB) initiative.
  \item Acted as the primary liaison between the organization and external stakeholders, successfully bridging communication with village officials and local communities to secure support for project implementation.
  \item Managed project documentation and timeline tracking to ensure all technical and operational milestones were achieved efficiently.
\end{itemize}

\resumeSubheading{Secretary of Minister of Information Technology — KM ITERA}{Apr 2024 – Dec 2024}{Lampung, Indonesia}{}
\vspace{-0.3em}
\begin{itemize}
  \item Managed meeting notes, project documentation, and Kanban-based project tracking through GitHub to coordinate IT-related student government web projects and report progress to the Minister of Information Technology
\end{itemize}

% \resumeSubheading{Head of Entrepreneurship Department — Sobat Bumi Lampung}{Oct 2025 – Present}{Lampung, Indonesia}{}
% \vspace{-0.3em}
% \begin{itemize}
%   \item Coordinated vendor communication and production planning for Sobat Bumi Lampung merchandise, including jackets, vests, and community attributes, while supporting entrepreneurship initiatives within the scholarship community
% \end{itemize}


% \resumeSubheading{Staff of Publication \& Documentation Division — HMIF ITERA}{Mar 2024 – Dec 2024}{Lampung, Indonesia}{}
% \begin{itemize}
%   \item Documented organizational events and managed public relations materials for the Informatics Student Association
% \end{itemize}

\section*{Skills}
\begin{itemize}
  \item \textbf{Business \& Leadership:} Cross-functional Collaboration, Stakeholder Management, Strategic Planning, Tech Community Leadership
  \item \textbf{Product \& Data Analysis:} Data-Driven Decision Making, Analytical Thinking, LLM-powered Applications, Product Ideation
  \item \textbf{Technical Core:} Python, SQL, Pandas, Machine Learning \& Deep Learning
  \item \textbf{Software \& Tools:} Advanced Microsoft Excel, Power BI, Notion, Trello, Figma (UI/UX Prototyping)
\end{itemize}

\section*{Certificats}
\begin{itemize}
  \item \textbf{ITERA English Proficiency Test (InciTe)} - Language Center ITERA (Apr 2026) | Score: 567
  \item \textbf{AWS re/Start Cloud Computing Program} - Orbit Future Academy \& Amazon Web Services (Mar 2026 - Present)
  \item \textbf{Belajar Pemrograman Prosedural dengan Python} - Dicoding (2024) | \textit{Credential ID: 2VX3O4W23ZYQ}
  \item \textbf{Memulai Pemrograman dengan Python} - Dicoding (2024) | \textit{Credential ID: L4PQQ8O2PO01}
  \item \textbf{Belajar Dasar Structured Query Language (SQL)} - Dicoding (2024) | \textit{Credential ID: NVP77M3WVPR0}
  \item \textbf{AI Praktis untuk Productivity} - Dicoding (2024) | \textit{Credential ID: 1OP82QOYLPQK}
\end{itemize}

\section*{Awards \& Honors}
\begin{itemize}
  \item \textbf{Pertamina Sobat Bumi Scholarship Recipient} - Pertamina Foundation (2024)
  \item \textbf{2nd Winner of UI/UX Competition Point Project 2.0} - HMIF ITERA (2024) | \textit{Credential ID: 313/525/SRT/HMIF/XII/2024}
  \item Guest Speaker for Scholarship Sharing Session at IFGTPB - HMIF ITERA (2024)
\end{itemize}

% \section*{Languages \& Publications}
% \begin{itemize}
%   \item \textbf{Indonesian:} Native 
%   \item \textbf{English:} Professional Working Proficiency
%   \item \textbf{Publications:} Sianturi, E. E. Y., et al. \textit{"Transformasi Digital Layanan Posyandu."} PADIMAS, 2025 $|$ \href{https://jurnal.mdp.ac.id/index.php/padimas/article/view/12206}{\footnotesize [Full Paper]}
% \end{itemize}

\section*{Publications}
\begin{itemize}
  \item Sianturi, E. E. Y., et al. \textit{"Transformasi Digital Layanan Posyandu."} PADIMAS, 2025 $|$ \href{https://jurnal.mdp.ac.id/index.php/padimas/article/view/12206}{\footnotesize [Full Paper]}
\end{itemize}
